
\chapter{Introduction}


\section{Defining function}
The main unifying theme in this dissertation is gene, or protein, function.  While, on the surface this might seem like a simple concept, in reality it has multiple meanings and a tangled history within the biological sciences; with an even longer history within philosophy itself.  It is  prudent, as a starting point for this document to give a clear working example of which definition of function we are using.

\subsection{The history of gene function}

\subsection{Functional versus function}

\subsection{The acorn and the oak tree}

I should be clear here that we are not working with a teleological definition of function.   yet wish to keep my definition as pragmatic and grounded in epistemological observation.  




While it may be tempting to frame the concept of protein function
in teleological terms, in reality it is pragmatically defined by the
distinct and diverse experimental tools available to a biologist.
For example, experimentalists can detect how a protein affects the
rate at which a reaction catalyzes through enzymatic assays. Yeast
2-hybrid assays can indicate binding between pairs of proteins. Microarrays
can be used to approximate the expression levels of a gene in a tissue
sample. Results from different experiments are then brought together
in order to paint a picture of how a protein behaves, usually culminating
in bestowing upon the protein a particular function, most commonly,
in the form of a Gene Ontology (GO) term annotation 
%\cite{ashburner2000gene}.
It is both the widespread application of GO terms, and its standardization
of terminology, that has made much of the quantitative analysis of
protein function possible.




